\documentclass[10pt,a4paper]{amsart}
\usepackage{amsmath,amssymb,amsthm,graphicx}
\usepackage{verbatim}
\newenvironment{reference}{\comment}{\endcomment}
%\newenvironment{reference}{}{}
\newenvironment{slogan}{\comment}{\endcomment}
\newenvironment{history}{\comment}{\endcomment}

% For commutative diagrams we use Xy-pic
\usepackage[all]{xy}

% We use 2cell for 2-commutative diagrams.
\xyoption{2cell}
% Theorem environments.
%
\theoremstyle{plain}
\newtheorem{theorem}[subsection]{Theorem}
\newtheorem{proposition}[subsection]{Proposition}
\newtheorem{lemma}[subsection]{Lemma}

\theoremstyle{definition}
\newtheorem{definition}[subsection]{Definition}
\newtheorem{example}[subsection]{Example}
\newtheorem{exercise}[subsection]{Exercise}
\newtheorem{situation}[subsection]{Situation}

\theoremstyle{remark}
\newtheorem{remark}[subsection]{Remark}
\newtheorem{remarks}[subsection]{Remarks}

\numberwithin{equation}{subsection}

% Macros
%
\def\lim{\mathop{\mathrm{lim}}\nolimits}
\def\colim{\mathop{\mathrm{colim}}\nolimits}
\def\Spec{\mathop{\mathrm{Spec}}}
\def\Hom{\mathop{\mathrm{Hom}}\nolimits}
\def\Ext{\mathop{\mathrm{Ext}}\nolimits}
\def\SheafHom{\mathop{\mathcal{H}\!\mathit{om}}\nolimits}
\def\SheafExt{\mathop{\mathcal{E}\!\mathit{xt}}\nolimits}
\def\Sch{\mathit{Sch}}
\def\Mor{\mathop{\mathrm{Mor}}\nolimits}
\def\Ob{\mathop{\mathrm{Ob}}\nolimits}
\def\Sh{\mathop{\mathit{Sh}}\nolimits}
\def\NL{\mathop{N\!L}\nolimits}
\def\CH{\mathop{\mathrm{CH}}\nolimits}
\def\proetale{{pro\text{-}\acute{e}tale}}
\def\etale{{\acute{e}tale}}
\def\QCoh{\mathit{QCoh}}
\def\Ker{\mathop{\mathrm{Ker}}}
\def\Im{\mathop{\mathrm{Im}}}
\def\Coker{\mathop{\mathrm{Coker}}}
\def\Coim{\mathop{\mathrm{Coim}}}

% Boxtimes
%
\DeclareMathSymbol{\boxtimes}{\mathbin}{AMSa}{"02}

%
% Macros for moduli stacks/spaces
%
\def\QCohstack{\mathcal{QC}\!\mathit{oh}}
\def\Cohstack{\mathcal{C}\!\mathit{oh}}
\def\Spacesstack{\mathcal{S}\!\mathit{paces}}
\def\Quotfunctor{\mathrm{Quot}}
\def\Hilbfunctor{\mathrm{Hilb}}
\def\Curvesstack{\mathcal{C}\!\mathit{urves}}
\def\Polarizedstack{\mathcal{P}\!\mathit{olarized}}
\def\Complexesstack{\mathcal{C}\!\mathit{omplexes}}
% \Pic is the operator that assigns to X its picard group, usage \Pic(X)
% \Picardstack_{X/B} denotes the Picard stack of X over B
% \Picardfunctor_{X/B} denotes the Picard functor of X over B
\def\Pic{\mathop{\mathrm{Pic}}\nolimits}
\def\Picardstack{\mathcal{P}\!\mathit{ic}}
\def\Picardfunctor{\mathrm{Pic}}
\def\Deformationcategory{\mathcal{D}\!\mathit{ef}}
\usepackage{cleveref}
\crefname{originallemma}{original lemma}{original lemmas}
\Crefname{originallemma}{Original lemma}{Original lemmas}
\makeatletter
\newlabel{categories-lemma-left-dual}{{0FFQ}{source}{}{sourcepointer-1}{}}
\newlabel{categories-lemma-left-dual@cref}{{[originallemma][0][]0FFQ}{[1][0][]source}{}{}{}}
\makeatother
\crefname{originalsection}{original section}{original sections}
\Crefname{originalsection}{Original section}{Original sections}
\makeatletter
\newlabel{section-karoubian}{{09SF}{source}{}{sourcepointer-2}{}}
\newlabel{section-karoubian@cref}{{[originalsection][0][]09SF}{[1][0][]source}{}{}{}}
\makeatother
\crefname{originalremark}{original remark}{original remarks}
\Crefname{originalremark}{Original remark}{Original remarks}
\makeatletter
\newlabel{categories-remark-left-dual-adjoint}{{0FFR}{source}{}{sourcepointer-3}{}}
\newlabel{categories-remark-left-dual-adjoint@cref}{{[originalremark][0][]0FFR}{[1][0][]source}{}{}{}}
\makeatother
\title{Every summand of an object with a left dual has a left dual}
\author{The Stacks Project Authors}
\date{stacks-project:homology.tex:local-snapshot-20261012; GNU FDL 1.2 or later}
\begin{document}
\maketitle
\noindent\textit{Editorial note.} All delivered mathematical text is the author's own, copied from the pinned Stacks Project chapter file homology.tex; no statement or proof step was rewritten, reordered or removed.\par\medskip
\noindent\textit{Editorial note.} The proof is the author's complete argument; the opening sentence 'We will use Categories, Lemma categories-lemma-left-dual without further mention.' names the same technique applied at each step rather than declining a step, and every identification the argument rests on is displayed or stated inside the counted span.\par\medskip
\noindent\textit{Editorial note.} Every labelled reference occurring inside the counted spans is carried in assembly.external\_references as a non-load-bearing navigation pointer (same-chapter ones at their definition site, cross-chapter ones at the invocation inside the counted span), so no load-bearing mathematics is replaced by a stub. No bibliography citation occurs inside the counted spans.\par\medskip
\noindent\textit{Editorial note.} The same-name check for this card scanned every staged Stacks chapter of this repository -- all sources/stacks-*/pinned-original-source.tex files, the card's own chapter excluded -- and the labels found are recorded in dedup.same\_name\_cross\_chapter on the card.\par\medskip
\noindent\textit{Editorial note.} The preamble is the chapter's own native preamble excerpt already used for this source: the comment-type reference, slogan and history environments (native preamble lines 8-13), the Xy-pic setup (native preamble lines 14-18) and the native theorem declarations with the mathematical macros (native preamble lines 154-217, the actual last line of the file).\par\medskip
\noindent\textit{Allowed standard prerequisites (editorial).}\begin{itemize}
\item Left duals and the adjunction with the tensor product. Categories, Lemma categories-lemma-left-dual and Remark categories-remark-left-dual-adjoint. Supply the identification of Hom(X tensor Y, Z) with Hom(X, Z tensor Y) and the criterion for an object to be a left dual of another, which the proof uses throughout.
\item Karoubian categories and split idempotents. Same chapter, Section section-karoubian. Supplies that an idempotent endomorphism of an object produces a direct sum decomposition, applied to the idempotents obtained in the endomorphism algebra of the dual.
\end{itemize}\medskip
\ifdefined\hypertarget\hypertarget{sourcepointer-1}{}\fi
\noindent\textit{Source reference.} Original lemma 0FFQ: sources/stacks-categories/pinned-original-source.tex, source lines 9558--9558. Navigation pointer invoked at proof line 4223; the counted span contains the pointer only and the cited result is not reproduced inside it, so the pointer is registered as a non-load-bearing navigation stub in the established form for Stacks material. No bibliography citation is involved.\par\medskip
\ifdefined\hypertarget\hypertarget{sourcepointer-2}{}\fi
\noindent\textit{Source reference.} Original section 09SF: sources/stacks-homology/pinned-original-source.tex, source lines 725--725. Navigation pointer invoked at proof line 4240; the counted span contains the pointer only and the cited result is not reproduced inside it, so the pointer is registered as a non-load-bearing navigation stub in the established form for Stacks material. No bibliography citation is involved.\par\medskip
\ifdefined\hypertarget\hypertarget{sourcepointer-3}{}\fi
\noindent\textit{Source reference.} Original remark 0FFR: sources/stacks-categories/pinned-original-source.tex, source lines 9603--9603. Navigation pointer invoked at proof line 4213; the counted span contains the pointer only and the cited result is not reproduced inside it, so the pointer is registered as a non-load-bearing navigation stub in the established form for Stacks material. No bibliography citation is involved.\par\medskip
\begin{lemma}
\label{lemma-Karoubian-dual}
In a Karoubian additive monoidal category every summand
of an object which has a left dual has a left dual.
\end{lemma}

\begin{proof}
We will use Categories, Lemma \ref{categories-lemma-left-dual}
without further mention.
Let $X$ be an object which has a left dual $Y$. We have
$$
\Hom(X, X) = \Hom(\mathbf{1}, X \otimes Y) = \Hom(Y, Y)
$$
If $a : X \to X$ corresponds to $b : Y \to Y$ then $b$ is the unique
endomorphism of $Y$ such that precomposing by $a$ on
$$
\Hom(Z' \otimes X, Z) = \Hom(Z', Z \otimes Y)
$$
is the same as postcomposing by $1 \otimes b$.
Hence the bijection $\Hom(X, X) \to \Hom(Y, Y)$, $a \mapsto b$
is an isomorphism of the opposite of the algebra $\Hom(X, X)$ with
the algebra $\Hom(Y, Y)$. In particular, if $X = X_1 \oplus X_2$, then
the corresponding projectors $e_1, e_2$ are mapped to idempotents
in $\Hom(Y, Y)$. If $Y = Y_1 \oplus Y_2$ is the corresponding direct
sum decomposition of $Y$ (Section \ref{section-karoubian})
then we see that under the bijection
$\Hom(Z' \otimes X, Z) = \Hom(Z', Z \otimes Y)$
we have $\Hom(Z' \otimes X_i, Z) = \Hom(Z', Z \otimes Y_i)$
functorially as subgroups for $i = 1, 2$.
It follows that $Y_i$ is the left dual of
$X_i$ by the discussion in
Categories, Remark \ref{categories-remark-left-dual-adjoint}.
\end{proof}

\end{document}
